% Section VI: Quantization on surfaces and emergence of TQFT structure
% Purpose: Show how canonical quantization on M = ℝ × Σ leads to functorial structure
%
% Key idea: Phase space = moduli space of flat connections on Σ
%          Quantization → Hilbert space H(Σ) assigned to surface
%          This is the seed of functorial TQFT structure
%
% Status: EXTRACT AND REORGANIZE
% Sources:
% - writings/brian_chern_simons_theory.tex (canonical quantization on surfaces)
%   Best mathematical treatment of symplectic quotient and geometric quantization
% - tex_docs/dense_derivation_expansion.tex (torus Hilbert space)
%   Lines ~800-1200: Explicit calculation for T² → dim H(T²) = k
%   This is the key worked example
% - chern_simons_two_review_synthesis.tex (overview)
%   Connects to functorial picture
%
% Action needed: REORGANIZE to emphasize functorial emergence
% Make explicit: "This procedure assigns Hilbert spaces to surfaces and amplitudes to 3-manifolds"
% Prepare reader for formal TQFT definition in Section VII

\section{Canonical Quantization and Functorial Structure}

% TODO: Organize to show emergence of bordism functor structure

\subsection{Phase Space on a Spatial Slice}

% TODO: For M = ℝ × Σ, identify phase space
% Constraint F = 0 → moduli space M_flat(Σ,G)
% Symplectic structure from CS action
% Source: brian_chern_simons_theory.tex (detailed treatment)

\subsection{Geometric Quantization of Flat Connection Moduli}

% TODO: Quantization of M_flat(Σ,G) → Hilbert space H(Σ)
% Mention Kähler polarization, prequantum line bundle
% Keep this BRIEF - physics paper, not pure math
% Source: brian_chern_simons_theory.tex (simplify and shorten)

\subsection{Example: The Torus}

% TODO: Full worked example T² → H(T²)
% Show dim H(T²) = k explicitly
% Source: dense_derivation_expansion.tex, lines ~800-1200
% This is concrete and crucial

% FIGURE FLAG [F-torus-hilbert]: Torus with flat connections,
% moduli space T² parametrized by holonomies, k quantum states

\subsection{Functorial Perspective}

% TODO: Make explicit what we've done:
% - Surfaces Σ → Hilbert spaces H(Σ)
% - 3-manifolds M with boundary → amplitudes Z(M): H(∂M) → ℂ
% - This is a functor structure (formalized in Section VII)
%
% Include composition: gluing 3-manifolds along boundaries
% Path integral respects gluing → functoriality

% This section explicitly prepares for Atiyah-Segal definition coming next
